% ============================================================================
%  ANEXO A — MATERIAL COMPLEMENTARIO
% ============================================================================

\chapter{Material complementario}
\label{anexo:complementario}

Este anexo reúne antecedentes de respaldo para la caracterización funcional
presentada en los Capítulos~\ref{cap:metodologia} y~\ref{cap:resultados}. Se
incluyen el diccionario de variables institucionales, el listado de
establecimientos incluidos y excluidos, la asignación final de clústeres y los
principales parámetros del pipeline computacional.

\section{Diccionario de variables institucionales}
\label{anexo:diccionario}

La Tabla~\ref{tab:anexo-diccionario} detalla las 28 variables
clínico-operativas directas construidas a nivel institucional. La variable
\texttt{peso\_medio\_cma} se conserva para fines descriptivos, pero fue excluida
de la matriz empleada en el clustering debido a que no es estructuralmente
comparable entre todos los establecimientos. Las 27 variables directas restantes,
junto con las 8 componentes principales de casuística
(\texttt{dim\_01} a \texttt{dim\_08}), conforman la matriz final de
$65 \times 35$ variables descrita en la Sección~\ref{sec:matriz-final}.

\begin{longtable}{p{4.0cm} p{5.8cm} p{2.0cm} p{2.2cm}}
\caption{Diccionario de las 28 variables clínico-operativas directas construidas a nivel institucional.}
\label{tab:anexo-diccionario} \\
\toprule
\textbf{Variable} & \textbf{Definición operacional} & \textbf{Unidad} & \textbf{Grupo} \\
\midrule
\endfirsthead

\multicolumn{4}{c}{\tablename\ \thetable{} -- continuación} \\
\toprule
\textbf{Variable} & \textbf{Definición operacional} & \textbf{Unidad} & \textbf{Grupo} \\
\midrule
\endhead

\texttt{egresos\_por\_anio} &
Promedio anual de egresos de hospitalización convencional durante el período analizado. &
Conteo & Actividad \\

\texttt{estancia\_media} &
Media aritmética de los días de estancia entre ingreso y alta. &
Días & Estadía \\

\texttt{estancia\_mediana} &
Mediana de los días de estancia entre ingreso y alta. &
Días & Estadía \\

\texttt{peso\_medio\_grd} &
Promedio del peso relativo GRD de los egresos del establecimiento. &
Adimensional & Complejidad \\

\texttt{severidad\_media} &
Promedio del índice de severidad clínica asignado por el agrupador GRD. &
Adimensional (1--4) & Complejidad \\

\texttt{mortalidad\_media} &
Promedio del índice de riesgo de mortalidad asignado por el agrupador GRD. &
Adimensional (1--4) & Complejidad \\

\texttt{entropia\_grd} &
Entropía de Shannon de la distribución de GRD del establecimiento. &
Bits & Diversidad \\

\texttt{comorbilidades\_promedio} &
Promedio de diagnósticos secundarios registrados por egreso. &
Conteo & Diversidad \\

\texttt{tasa\_cma} &
Proporción de egresos registrados como Cirugía Mayor Ambulatoria sobre el total institucional. &
Proporción & CMA \\

\texttt{peso\_medio\_cma} &
Promedio del peso relativo GRD entre los egresos de CMA. Se emplea de forma descriptiva y se excluye del clustering final. &
Adimensional & CMA \\

\texttt{pct\_pediatrico} &
Proporción de egresos de pacientes menores de 15 años. &
Proporción & Perfil demográfico \\

\texttt{pct\_geriatrico} &
Proporción de egresos de pacientes de 65 años o más. &
Proporción & Perfil demográfico \\

\texttt{pct\_femenino\_fertil} &
Proporción de egresos de mujeres entre 15 y 49 años. &
Proporción & Perfil demográfico \\

\texttt{edad\_mediana} &
Mediana de edad al ingreso de los egresos institucionales. &
Años & Perfil demográfico \\

\texttt{pct\_urgencia} &
Proporción de egresos cuyo tipo de ingreso fue urgencia. &
Proporción & Modalidad de ingreso \\

\texttt{pct\_programada} &
Proporción de egresos cuyo tipo de ingreso fue programado. &
Proporción & Modalidad de ingreso \\

\texttt{pct\_obstetrica\_ingreso} &
Proporción de egresos cuyo tipo de ingreso fue obstétrico. &
Proporción & Modalidad de ingreso \\

\texttt{pct\_alta\_domicilio} &
Proporción de egresos con alta a domicilio. &
Proporción & Tipo de alta \\

\texttt{pct\_alta\_fallecido} &
Proporción de egresos con alta por defunción. &
Proporción & Tipo de alta \\

\texttt{pct\_alta\_traslado} &
Proporción de egresos con alta por traslado o derivación a otro establecimiento. &
Proporción & Tipo de alta \\

\texttt{pct\_origen\_emergencia} &
Proporción de egresos procedentes del servicio de urgencia propio. &
Proporción & Procedencia \\

\texttt{pct\_origen\_referencia} &
Proporción de egresos derivados desde otro establecimiento de la red. &
Proporción & Procedencia \\

\texttt{pct\_uso\_pabellon} &
Proporción de egresos con al menos un uso de pabellón quirúrgico registrado. &
Proporción & Actividad quirúrgica \\

\texttt{pabellones\_promedio} &
Promedio de usos de pabellón por egreso hospitalizado. &
Conteo & Actividad quirúrgica \\

\texttt{tasa\_partos} (nombre técnico legado) &
Proporción de egresos hospitalarios con condición de alta neonatal registrada; no equivale a una tasa ni a un conteo de partos atendidos. &
Proporción & Actividad neonatal \\

\texttt{tasa\_prematurez} &
Proporción de egresos neonatales registrados con peso de nacimiento inferior a 2.500 gramos. &
Proporción & Actividad neonatal \\

\texttt{pct\_estancia\_larga} &
Proporción de egresos con estancia superior a 30 días. &
Proporción & Estadía \\

\texttt{cv\_estancia} &
Coeficiente de variación de la duración de estancia. &
Adimensional & Estadía \\

\bottomrule
\end{longtable}

Las ocho componentes de casuística
(\texttt{dim\_01}--\texttt{dim\_08}) se obtuvieron mediante PCA de la matriz
de proporciones diagnósticas, de procedimientos y de GRD descrita en la
Sección~\ref{sec:pca-detalle}. Cada componente es una combinación lineal de las
variables originales, por lo que no posee una definición operacional única. Su
interpretación se realiza mediante las cargas de mayor magnitud y debe
considerarse complementaria a las variables institucionales directas.

\section{Establecimientos incluidos, excluidos y asignación final}
\label{anexo:hospitales}

De los 72 establecimientos con datos GRD disponibles durante 2019--2024, 65
cumplieron el criterio de presencia temporal mínima de tres de los seis años del
período y conformaron el universo final de análisis. Los siete establecimientos
excluidos se presentan en la Tabla~\ref{tab:anexo-excluidos}.

\begin{table}[htbp]
\centering
\caption{Establecimientos excluidos por no cumplir el criterio de presencia temporal mínima.}
\label{tab:anexo-excluidos}
\begin{tabular}{c c}
\toprule
\textbf{Código} & \textbf{Años con presencia válida} \\
\midrule
110110 & 2 \\
116107 & 1 \\
116111 & 1 \\
118106 & 2 \\
119101 & 2 \\
119102 & 1 \\
200717 & 1 \\
\bottomrule
\end{tabular}
\end{table}

La partición final de Nivel~1 contiene tres hospitales pediátricos ($C0$), un
núcleo generalista de 54 establecimientos ($C1$), un grupo de cinco hospitales
con carga clínica y estancia prolongada ($C2$), y tres institutos monográficos
($C3$). El núcleo generalista se subdividió de forma exploratoria mediante
Ward con $K_{sub}=10$, cuyas etiquetas se presentan como $S0$ a $S9$.

La Tabla~\ref{tab:anexo-hospitales} presenta la asignación final de los 65
establecimientos. Las etiquetas de Nivel~2 solo se aplican a los miembros de
$C1$; no constituyen categorías funcionales validadas equivalentes a los grupos
del Nivel~1.

\begin{longtable}{p{1.5cm} p{9.0cm} p{1.0cm} p{1.0cm}}
\caption{Asignación final de los establecimientos en los Niveles 1 y 2 del clustering Ward.}
\label{tab:anexo-hospitales} \\
\toprule
\textbf{Código} & \textbf{Establecimiento} & \textbf{N1} & \textbf{N2} \\
\midrule
\endfirsthead

\multicolumn{4}{c}{\tablename\ \thetable{} -- continuación} \\
\toprule
\textbf{Código} & \textbf{Establecimiento} & \textbf{N1} & \textbf{N2} \\
\midrule
\endhead

101100 & Hospital Dr.\ Juan Noé Crevanni (Arica) & $C1$ & $S6$ \\
102100 & Hospital Dr.\ Ernesto Torres Galdames (Iquique) & $C1$ & $S6$ \\
103100 & Hospital Dr.\ Leonardo Guzmán (Antofagasta) & $C1$ & $S8$ \\
103101 & Hospital Dr.\ Carlos Cisternas (Calama) & $C1$ & $S6$ \\
104100 & Hospital San José del Carmen (Copiapó) & $C1$ & $S6$ \\
104103 & Hospital Provincial del Huasco Monseñor Fernando Ariztía Ruiz (Vallenar) & $C1$ & $S6$ \\
105100 & Hospital San Juan de Dios (La Serena) & $C1$ & $S5$ \\
105101 & Hospital San Pablo (Coquimbo) & $C1$ & $S6$ \\
105102 & Hospital Dr.\ Antonio Tirado Lanas (Ovalle) & $C1$ & $S5$ \\
106100 & Hospital Carlos Van Buren (Valparaíso) & $C1$ & $S8$ \\
106102 & Hospital Dr.\ Eduardo Pereira Ramírez (Valparaíso) & $C3$ & --- \\
106103 & Hospital Claudio Vicuña (San Antonio) & $C1$ & $S0$ \\
107100 & Hospital Dr.\ Gustavo Fricke (Viña del Mar) & $C1$ & $S9$ \\
107101 & Hospital San Martín (Quillota) & $C1$ & $S5$ \\
107102 & Hospital de Quilpué & $C1$ & $S1$ \\
108100 & Hospital de San Camilo (San Felipe) & $C1$ & $S5$ \\
108101 & Hospital San Juan de Dios (Los Andes) & $C1$ & $S6$ \\
109100 & Complejo Hospitalario San José (Santiago, Independencia) & $C1$ & $S3$ \\
109101 & Hospital Clínico de Niños Dr.\ Roberto del Río & $C0$ & --- \\
110100 & Hospital San Juan de Dios (Santiago) & $C1$ & $S8$ \\
110120 & Hospital Dr.\ Félix Bulnes Cerda & $C1$ & $S3$ \\
110130 & Hospital Adalberto Steeger (Talagante) & $C1$ & $S2$ \\
110150 & Hospital San José (Melipilla) & $C1$ & $S2$ \\
111100 & Hospital Clínico San Borja-Arriarán & $C1$ & $S4$ \\
111101 & Hospital Clínico Metropolitano El Carmen Dr.\ Luis Valentín Ferrada & $C1$ & $S7$ \\
111195 & Hospital de Urgencia Asistencia Pública Dr.\ Alejandro del Río & $C2$ & --- \\
112100 & Hospital Del Salvador (Santiago, Providencia) & $C2$ & --- \\
112101 & Hospital Dr.\ Luis Tisné B. & $C1$ & $S3$ \\
112102 & Hospital de Niños Dr.\ Luis Calvo Mackenna & $C0$ & --- \\
112103 & Instituto Nacional de Enfermedades Respiratorias y Cirugía Torácica & $C3$ & --- \\
112104 & Instituto de Neurocirugía Dr.\ Alfonso Asenjo & $C3$ & --- \\
113100 & Hospital Barros Luco Trudeau & $C1$ & $S7$ \\
113130 & Hospital Dr.\ Exequiel González Cortés & $C0$ & --- \\
113150 & Hospital San Luis (Buin) & $C1$ & $S2$ \\
113180 & Hospital El Pino & $C1$ & $S3$ \\
114101 & Complejo Hospitalario Dr.\ Sótero del Río & $C1$ & $S8$ \\
114103 & Hospital Padre Alberto Hurtado & $C1$ & $S3$ \\
114105 & Hospital Clínico Metropolitano La Florida Dra.\ Eloísa Díaz Inzunza & $C1$ & $S7$ \\
115100 & Hospital Regional de Rancagua & $C1$ & $S7$ \\
115107 & Hospital San Juan de Dios (San Fernando) & $C1$ & $S6$ \\
115110 & Hospital de Santa Cruz & $C1$ & $S1$ \\
116100 & Hospital San Juan de Dios (Curicó) & $C1$ & $S5$ \\
116105 & Hospital Dr.\ César Garavagno Burotto (Talca) & $C1$ & $S8$ \\
116108 & Hospital Presidente Carlos Ibáñez del Campo (Linares) & $C1$ & $S6$ \\
116110 & Hospital San José (Parral) & $C1$ & $S0$ \\
117101 & Hospital Clínico Herminda Martín (Chillán) & $C1$ & $S7$ \\
117102 & Hospital de San Carlos & $C1$ & $S6$ \\
118100 & Hospital Clínico Regional Dr.\ Guillermo Grant Benavente (Concepción) & $C1$ & $S8$ \\
118105 & Hospital San José (Coronel) & $C1$ & $S0$ \\
119100 & Hospital Las Higueras (Talcahuano) & $C1$ & $S6$ \\
120101 & Complejo Asistencial Dr.\ Víctor Ríos Ruiz (Los Ángeles) & $C1$ & $S6$ \\
121109 & Hospital Dr.\ Hernán Henríquez Aravena (Temuco) & $C1$ & $S8$ \\
121110 & Hospital Dr.\ Abraham Godoy (Lautaro) & $C2$ & --- \\
121114 & Hospital Intercultural de Nueva Imperial & $C2$ & --- \\
121117 & Hospital de Pitrufquén & $C2$ & --- \\
121121 & Hospital de Villarrica & $C1$ & $S1$ \\
122100 & Hospital Clínico Regional (Valdivia) & $C1$ & $S8$ \\
123100 & Hospital Base San José de Osorno & $C1$ & $S6$ \\
124105 & Hospital de Puerto Montt & $C1$ & $S7$ \\
125100 & Hospital Regional (Coihaique) & $C1$ & $S6$ \\
126100 & Hospital Clínico de Magallanes Dr.\ Lautaro Navarro Avaria & $C1$ & $S8$ \\
128109 & Hospital Provincial Dr.\ Rafael Avaría (Curanilahue) & $C1$ & $S1$ \\
129100 & Hospital Dr.\ Mauricio Heyermann (Angol) & $C1$ & $S6$ \\
129106 & Hospital San José (Victoria) & $C1$ & $S6$ \\
133150 & Hospital de Castro & $C1$ & $S6$ \\

\bottomrule
\end{longtable}

La asignación completa se encuentra en el archivo
\texttt{reports/tables/asignacion\_jerarquica\_final.csv}. Las
comparaciones con la clasificación administrativa MINSAL se presentan en la
Sección~\ref{sec:res-minsal}; no se reproducen en esta tabla para concentrar el
anexo en las asignaciones finales del procedimiento de clustering.

\section{Parámetros del pipeline}
\label{anexo:parametros}

La Tabla~\ref{tab:anexo-parametros} consolida las decisiones metodológicas y
los parámetros principales de la ejecución final. Las decisiones se justifican
en detalle en el Capítulo~\ref{cap:metodologia}; aquí se presentan como
referencia de implementación y reproducibilidad.

\begin{longtable}{p{3.6cm} p{4.3cm} p{5.8cm}}
\caption{Parámetros y decisiones de diseño del pipeline computacional.}
\label{tab:anexo-parametros} \\
\toprule
\textbf{Etapa} & \textbf{Parámetro} & \textbf{Valor o decisión} \\
\midrule
\endfirsthead

\multicolumn{3}{c}{\tablename\ \thetable{} -- continuación} \\
\toprule
\textbf{Etapa} & \textbf{Parámetro} & \textbf{Valor o decisión} \\
\midrule
\endhead

Elegibilidad &
Presencia temporal mínima &
Al menos 3 de los 6 años del período 2019--2024. \\

Universo final &
Establecimientos incluidos &
65 hospitales públicos de alta y mediana complejidad. \\

Variables directas &
Variables construidas &
28 variables institucionales; \texttt{peso\_medio\_cma} se excluye de la matriz de clustering. \\

Matriz final &
Dimensión &
65 establecimientos $\times$ 35 variables: 27 variables directas y 8 componentes PCA. \\

Preprocesamiento &
Transformación de asimetría &
Transformación $\log(1+x)$ aplicada a variables con asimetría marcada, según la configuración de la ejecución final. \\

Preprocesamiento &
Escalado &
\textit{RobustScaler}, basado en mediana y rango intercuartílico. \\

Casuística &
Reducción dimensional &
PCA sobre la representación de casuística diagnóstica, de procedimientos y GRD. \\

Casuística &
Componentes retenidas &
8 componentes principales, con $72{,}5\,\%$ de varianza acumulada. \\

Clustering Nivel 1 &
Algoritmo &
Clustering jerárquico aglomerativo Ward con distancia euclidiana. \\

Clustering Nivel 1 &
Valores evaluados &
$K \in \{2,\ldots,10\}$. \\

Clustering Nivel 1 &
Solución reportada &
$K=4$, seleccionada como compromiso interpretativo; tamaños
$[3,54,5,3]$. \\

Clustering Nivel 2 &
Subconjunto analizado &
Núcleo generalista $C1$ ($n=54$). \\

Clustering Nivel 2 &
Algoritmo &
Clustering jerárquico aglomerativo Ward. \\

Clustering Nivel 2 &
Solución reportada &
$K_{sub}=10$, de carácter exploratorio. \\

Caracterización por variable &
Prueba &
Kruskal--Wallis con corrección de Benjamini--Hochberg. \\

Caracterización por variable &
Tamaño de efecto &
$\varepsilon^2$ para variables continuas. \\

Atipicidad &
Criterios &
Isolation Forest global, distancia al centroide normalizada por MAD global y Silhouette individual negativo. \\

Atipicidad &
Sensibilidad &
Isolation Forest evaluado con contaminación de 5\,\%, 10\,\%, 15\,\% y 20\,\%. \\

Sensibilidad algorítmica &
Métodos &
K-means ($K=4$), GMM ($K=4$) y MST-kNN para sensibilidad interna. \\

Sensibilidad algorítmica &
MST-kNN &
Evaluación con $k \in \{3,\ldots,8\}$; los resultados se interpretan como sensibilidad interna y no como validación externa. \\

Sensibilidad temporal &
Comparación &
Período completo 2019--2024 frente a 2019 y 2022--2024, excluyendo 2020--2021. \\

Sensibilidad composicional &
Representación alternativa &
Transformación CLR de la casuística antes de PCA. \\

Reproducibilidad &
Semilla aleatoria &
\texttt{RANDOM\_STATE = 42} en los algoritmos estocásticos que admiten semilla. \\

Reproducibilidad &
Implementación &
Python 3.11 y dependencias declaradas en \texttt{requirements.txt} y \texttt{pyproject.toml}. \\

\bottomrule
\end{longtable}

\section{Estructura del repositorio y reproducibilidad}
\label{anexo:repositorio}

El código del proyecto está disponible en
\url{https://github.com/reii23/chile-hospital-grd-profiling} y se organiza en
los siguientes directorios principales:

\begin{itemize}
  \item \texttt{data/raw/}: archivos fuente GRD organizados por año.

  \item \texttt{data/processed/}: matrices intermedias y finales generadas por
  el pipeline, almacenadas principalmente en formato Parquet.

  \item \texttt{src/etl/}: módulos de extracción, limpieza, homologación y
  construcción de variables.

  \item \texttt{src/modeling/}: módulos para construcción de matrices,
  reducción dimensional, clustering y comparación de métodos.

  \item \texttt{src/utils/}: funciones auxiliares de entrada, salida y
  validación de datos, con un registro único de rutas de artefactos.

  \item \texttt{scripts/}: ejecución no interactiva del pipeline. Comprende la
  construcción del insumo GRD y de la matriz institucional
  (\texttt{build\_*.py}), los contrastes estadísticos y los análisis de
  sensibilidad (\texttt{analyze\_*.py}, \texttt{compare\_*.py}) y la
  regeneración de las figuras incluidas en este documento
  (\texttt{regenerar\_figuras\_*.py}).

  \item \texttt{notebooks/}: registro exploratorio del proceso analítico.

  \item \texttt{tests/}: pruebas unitarias y basadas en propiedades sobre los
  módulos de \texttt{src/}, ejecutadas con \texttt{pytest} y
  \texttt{hypothesis}.

  \item \texttt{reports/tables/}: resultados tabulares en formato CSV. Todas
  las cifras reportadas en el Capítulo~\ref{cap:resultados} provienen de estos
  archivos.

  \item \texttt{reports/figures/}: figuras exploratorias auxiliares. Las
  figuras incluidas en este documento se generan en
  \texttt{formato-tesis/}, junto con las fuentes \LaTeX{} del informe.

  \item \texttt{requirements.txt}, \texttt{requirements-dev.txt} y
  \texttt{pyproject.toml}: declaración de dependencias del entorno de ejecución.
\end{itemize}

Los insumos necesarios para ejecutar el pipeline se distribuyen junto con el
código: las tablas maestras de clasificación clínica (CIE-10, CIE-9-MC y la
maestra de establecimientos de las bases GRD), la Base de Establecimientos del
DEIS empleada para recuperar el Nivel de Complejidad asignado por el MINSAL, y
los artefactos intermedios en formato Parquet, entre ellos la matriz
institucional de 65 hospitales y 38 columnas sobre la que operan los análisis.

Se exceptúan dos artefactos que exceden el límite de tamaño por archivo de la
plataforma de alojamiento: los seis archivos anuales de egresos GRD
($3{,}8$\,GB en total), que se descargan del Tablero GRD del
MINSAL~\citep{TableauGRD}, y el archivo consolidado
\texttt{grd\_filtrado.parquet} ($198$\,MB), que se regenera desde
los anteriores mediante \texttt{build\_extended\_grd.py}. Con esa excepción, 15
de los 21 scripts se ejecutan directamente sobre una copia del repositorio; los
seis restantes recorren los egresos individuales y requieren la descarga previa.

El código se publica bajo licencia MIT, de modo que puede reutilizarse,
modificarse e integrarse en trabajos posteriores conservando la atribución de
autoría. Esta licencia cubre únicamente el código: los egresos GRD, las
clasificaciones CIE y la Base de Establecimientos del DEIS se rigen por las
condiciones de sus titulares originales. El documento de esta tesis se publica
por separado bajo Licencia Creative Commons Atribución-Chile 3.0.

El orden de ejecución es relevante, pues los análisis de sensibilidad y los
contrastes estadísticos consumen la partición reportada. Esta se fija en
\texttt{build\_hierarchical\_clustering.py} y se persiste en
\texttt{reports/tables/asignacion\_jerarquica\_final.csv}. El archivo
\texttt{README.md} del repositorio detalla la secuencia completa. Junto con los
parámetros consolidados en la Tabla~\ref{tab:anexo-parametros} y las semillas
aleatorias fijadas en el código, esta organización permite reconstruir el flujo
de procesamiento y los análisis descritos en el
Capítulo~\ref{cap:metodologia}.

\section{Detalle suplementario del Nivel 2}
\label{anexo:nivel2-exploratorio}

La solución Ward de $K_{sub}=10$ sobre los 54 hospitales del núcleo generalista
se conserva únicamente como exploración suplementaria. La validación por
submuestreo mostró una concordancia mediana de ARI $=0{,}476$ y NMI $=0{,}748$,
insuficiente para sostener una tipología funcional estable. Por ello, las
etiquetas $S0$--$S9$ de la Tabla~\ref{tab:anexo-hospitales} no son categorías
validadas ni deben utilizarse para inferencia. En particular, $S4$ (Hospital
Clínico San Borja--Arriarán) y $S9$ (Hospital Dr. Gustavo Fricke) son
observaciones unitarias y se interpretan como establecimientos alejados del
centroide del núcleo, no como subtipos colectivos.

Los resultados reproducibles de las 100 réplicas, matrices de coasignación y
estabilidad individual se entregan en los archivos
\texttt{reports/tables/estabilidad\_nivel2\_*.csv} del repositorio.
